\documentclass[11pt,a4paper]{article}
\usepackage[utf8]{inputenc}
\usepackage[T1]{fontenc}
\usepackage[spanish]{babel}
\usepackage{amsmath,amssymb}
\usepackage{booktabs}
\usepackage{longtable}
\usepackage{enumitem}
\usepackage{geometry}
\usepackage{graphicx}
\usepackage{hyperref}
\usepackage{listings}
\usepackage{xcolor}

\geometry{margin=2.5cm}

\lstset{
  basicstyle=\ttfamily\small,
  breaklines=true,
  frame=single,
  numbers=left,
  numberstyle=\tiny\color{gray},
  keywordstyle=\color{blue},
  commentstyle=\color{green!50!black},
  stringstyle=\color{red!70!black},
}

\title{%
  \textbf{Informe del Proyecto}\\[6pt]
  \large Parsers LL(1) y SLR(1) con Evaluador de SDD\\[4pt]
  \normalsize SI2002 / ST0270 / CM0081 --- Lenguajes Formales\\
  Universidad EAFIT, 2026-2\\
  Profesor: Oscar Rodr\'iguez Cifuentes
}
\author{
  Samuel Gonz\'alez P\'erez\\
  \texttt{sgonzalep3@eafit.edu.co}
  \and
  Samuel L\'opez Bonilla\\
  \texttt{slopezb8q@eafit.edu.co}
}
\date{Noviembre 2026}

\begin{document}
\maketitle

%% ============================================================
\section{Algoritmos y estructuras de datos}
\label{sec:algorithms}

Todos los algoritmos siguen \emph{Compilers: Principles, Techniques, and Tools} (Aho et al., 2da edici\'on). La Tabla~\ref{tab:refs} lista las referencias espec\'ificas.

\begin{table}[h]
\centering
\caption{Referencias del libro para cada componente.}
\label{tab:refs}
\begin{tabular}{@{}lll@{}}
\toprule
\textbf{Componente} & \textbf{Secci\'on / Algoritmo} & \textbf{Estructura de datos} \\
\midrule
FIRST / FOLLOW          & Secci\'on 4.4.2                       & \texttt{dict[str, set[str]]} \\
Tabla LL(1)             & Secci\'on 4.4.3                       & \texttt{dict[(str,str), (str,list)]} \\
Parser predictivo       & Algoritmo 4.34, Secci\'on 4.4.4       & Pila expl\'icita (list) \\
\'Items LR(0)            & Secci\'on 4.6.2                       & Clase \texttt{Item} (head, body, dot) \\
CLOSURE / GOTO          & Algoritmo 4.32, Secci\'on 4.6.2       & \texttt{frozenset[Item]} \\
Colecci\'on can\'onica    & Secci\'on 4.6.2                       & Lista de conjuntos de \'items \\
Tablas SLR(1)           & Algoritmo 4.46, Secci\'on 4.6.4       & ACTION: \texttt{dict[(int,str), tuple]}, GOTO: \texttt{dict[(int,str), int]} \\
Parser LR               & Algoritmo 4.44, Secci\'on 4.6.3       & Pila de estados + pila de s\'imbolos \\
Grafo de dependencias   & Secci\'on 5.2.1                       & Lista de adyacencia + mapa de in-degree \\
Orden topol\'ogico       & Secci\'on 5.2.2 (algoritmo de Kahn)   & Cola (ordenada para determinismo) \\
Evaluador SDD           & Secciones 5.1--5.2                   & \'Arbol de parseo (\texttt{TreeNode}) \\
\bottomrule
\end{tabular}
\end{table}

La gram\'atica se representa como una clase de Python con conjuntos de no-terminales y terminales, un diccionario de producciones y un s\'imbolo inicial. Las producciones mapean cada no-terminal a una lista de alternativas, donde cada alternativa es una lista de s\'imbolos de un solo car\'acter. Epsilon se representa con el car\'acter~\texttt{e} y el marcador de fin de entrada con~\texttt{\$}.


%% ============================================================
\section{\'Items 1 y 2 --- FIRST/FOLLOW y Parsers}
\label{sec:items12}

\subsection{Gram\'atica propia}

Definimos la siguiente gram\'atica $G$:

\begin{align*}
S &\to a\,B\,S \mid b \\
B &\to b \mid \varepsilon
\end{align*}

En formato de entrada:
\begin{lstlisting}
2
S -> aBS b
B -> b e
\end{lstlisting}

$S$ es recursiva ($S \to aBS$ contiene $S$ en el cuerpo) y $B$ tiene una producci\'on-$\varepsilon$, as\'i que la gram\'atica ejercita los casos principales de FIRST/FOLLOW.

\subsection{Conjuntos FIRST y FOLLOW}

Para la gram\'atica $G$:

\begin{center}
\begin{tabular}{@{}lll@{}}
\toprule
\textbf{S\'imbolo} & \textbf{FIRST} & \textbf{FOLLOW} \\
\midrule
$S$ & $\{a, b\}$     & $\{\$\}$   \\
$B$ & $\{b, \varepsilon\}$    & $\{a, b\}$ \\
\bottomrule
\end{tabular}
\end{center}

Calculados por nuestro programa:
\begin{lstlisting}
First(B) = {b, e}
First(S) = {a, b}
Follow(B) = {a, b}
Follow(S) = {$}
\end{lstlisting}

\subsection{Tabla LL(1)}

\begin{center}
\begin{tabular}{@{}l|ccc@{}}
\toprule
 & $a$ & $b$ & $\$$ \\
\midrule
$S$ & $S \to aBS$ & $S \to b$ & \\
$B$ & $B \to \varepsilon$ & $B \to b$ & \\
\bottomrule
\end{tabular}
\end{center}

Ninguna celda tiene dos entradas, por lo tanto $G$ es LL(1).

\subsection{Parseo de \texorpdfstring{$w = abb$}{w = abb} paso a paso (LL(1))}

\begin{center}
\begin{longtable}{@{}clll@{}}
\toprule
\textbf{Paso} & \textbf{Pila} & \textbf{Entrada} & \textbf{Acci\'on} \\
\midrule
\endhead
1 & $\$\;S$       & $a\,b\,b\,\$$ & $S \to aBS$ \\
2 & $\$\;S\,B\,a$ & $a\,b\,b\,\$$ & match $a$ \\
3 & $\$\;S\,B$    & $b\,b\,\$$ & $B \to b$ \\
4 & $\$\;S\,b$    & $b\,b\,\$$ & match $b$ \\
5 & $\$\;S$       & $b\,\$$ & $S \to b$ \\
6 & $\$\;b$       & $b\,\$$ & match $b$ \\
7 & $\$$          & $\$$ & \textbf{aceptar} \\
\bottomrule
\end{longtable}
\end{center}

La cadena $abb$ se acepta en 7 pasos.


%% ============================================================
\section{\'Item 3 --- Definici\'on dirigida por sintaxis}
\label{sec:sdd}

\subsection{Paso 1: Gram\'atica}

Definimos un mini-calculador donde cada sentencia asigna una expresi\'on aritm\'etica a la variable~$x$:

\begin{align*}
S &\to x\,{=}\,E        & F &\to (\,E\,) \\
E &\to T\,A             & F &\to 1 \\
A &\to {+}\,T\,A \mid \varepsilon  & F &\to 2 \\
T &\to F\,M             & F &\to 3 \\
M &\to {*}\,F\,M \mid \varepsilon  &   &
\end{align*}

6 no-terminales, 11 producciones, terminales $\{x, =, +, *, (, ), 1, 2, 3\}$.

En formato de entrada:
\begin{lstlisting}
6
S -> x=E
E -> TA
A -> +TA e
T -> FM
M -> *FM e
F -> (E) 1 2 3
\end{lstlisting}

Esta gram\'atica es LL(1) (verificado por nuestro programa) y usa la t\'ecnica est\'andar de factorizar $E/T$ en colas $A/M$ para eliminar la recursi\'on izquierda.

\subsection{Paso 2: Atributos y reglas sem\'anticas}

\subsubsection*{Tabla de atributos}

\begin{center}
\begin{tabular}{@{}llll@{}}
\toprule
\textbf{S\'imbolo} & \textbf{Atributo} & \textbf{Tipo} & \textbf{Significado} \\
\midrule
$S$ & \texttt{val} & sintetizado & valor de la expresi\'on \\
$S$ & \texttt{ops} & sintetizado & n\'umero de operaciones aritm\'eticas \\
$E$ & \texttt{val} & sintetizado & valor \\
$E$ & \texttt{ops} & sintetizado & conteo de operaciones \\
$A$ & \texttt{inh} & heredado    & valor acumulado desde la izquierda \\
$A$ & \texttt{val} & sintetizado & valor despu\'es de todas las sumas \\
$A$ & \texttt{ops} & sintetizado & conteo de operaciones \\
$T$ & \texttt{val} & sintetizado & valor \\
$T$ & \texttt{ops} & sintetizado & conteo de operaciones \\
$M$ & \texttt{inh} & heredado    & valor acumulado desde la izquierda \\
$M$ & \texttt{val} & sintetizado & valor despu\'es de todas las multiplicaciones \\
$M$ & \texttt{ops} & sintetizado & conteo de operaciones \\
$F$ & \texttt{val} & sintetizado & valor del factor \\
$F$ & \texttt{ops} & sintetizado & conteo de operaciones \\
\bottomrule
\end{tabular}
\end{center}

\subsubsection*{Reglas sem\'anticas}

\begin{center}
\begin{longtable}{@{}l l@{}}
\toprule
\textbf{Producci\'on} & \textbf{Reglas sem\'anticas} \\
\midrule
\endhead
$S \to x\,{=}\,E$ &
  $S.\texttt{val} = E.\texttt{val}$; \quad $S.\texttt{ops} = E.\texttt{ops}$ \\[6pt]

$E \to T\,A$ &
  $A.\texttt{inh} = T.\texttt{val}$; \quad $E.\texttt{val} = A.\texttt{val}$; \\
  & $E.\texttt{ops} = T.\texttt{ops} + A.\texttt{ops}$ \\[6pt]

$A \to {+}\,T\,A_1$ &
  $A_1.\texttt{inh} = A.\texttt{inh} + T.\texttt{val}$; \quad $A.\texttt{val} = A_1.\texttt{val}$; \\
  & $A.\texttt{ops} = 1 + T.\texttt{ops} + A_1.\texttt{ops}$ \\[6pt]

$A \to \varepsilon$ &
  $A.\texttt{val} = A.\texttt{inh}$; \quad $A.\texttt{ops} = 0$ \\[6pt]

$T \to F\,M$ &
  $M.\texttt{inh} = F.\texttt{val}$; \quad $T.\texttt{val} = M.\texttt{val}$; \\
  & $T.\texttt{ops} = F.\texttt{ops} + M.\texttt{ops}$ \\[6pt]

$M \to {*}\,F\,M_1$ &
  $M_1.\texttt{inh} = M.\texttt{inh} \times F.\texttt{val}$; \quad $M.\texttt{val} = M_1.\texttt{val}$; \\
  & $M.\texttt{ops} = 1 + F.\texttt{ops} + M_1.\texttt{ops}$ \\[6pt]

$M \to \varepsilon$ &
  $M.\texttt{val} = M.\texttt{inh}$; \quad $M.\texttt{ops} = 0$ \\[6pt]

$F \to (\,E\,)$ &
  $F.\texttt{val} = E.\texttt{val}$; \quad $F.\texttt{ops} = E.\texttt{ops}$ \\[6pt]

$F \to 1$ &
  $F.\texttt{val} = 1$; \quad $F.\texttt{ops} = 0$ \\[3pt]

$F \to 2$ &
  $F.\texttt{val} = 2$; \quad $F.\texttt{ops} = 0$ \\[3pt]

$F \to 3$ &
  $F.\texttt{val} = 3$; \quad $F.\texttt{ops} = 0$ \\
\bottomrule
\end{longtable}
\end{center}

\subsection{Paso 3: Clasificaci\'on}

La SDD es \textbf{L-atribuida} pero \textbf{no S-atribuida}.

\begin{itemize}
  \item \textbf{No es S-atribuida}: los atributos $A.\texttt{inh}$ y $M.\texttt{inh}$ son heredados, por lo que la SDD contiene atributos heredados y no puede ser S-atribuida.
  \item \textbf{Es L-atribuida}: cada atributo heredado depende solo de:
    \begin{enumerate}
      \item atributos heredados de la cabeza (ej.\ $A_1.\texttt{inh}$ depende de $A.\texttt{inh}$), o
      \item atributos de s\'imbolos a la izquierda en el cuerpo de la producci\'on (ej.\ $A.\texttt{inh}$ en $E \to TA$ depende de $T.\texttt{val}$, y $T$ est\'a a la izquierda de~$A$).
    \end{enumerate}
    Ning\'un atributo heredado depende de un s\'imbolo a su derecha en el cuerpo. Esto satisface la Definici\'on~5.4 del libro.
\end{itemize}

\subsection{Paso 4: Ejemplo a mano}

\textbf{Cadena de entrada}: \texttt{x=2*3+1}

\subsubsection*{\'Arbol de parseo (anotado)}

Salida de nuestro programa para \texttt{x=2*3+1}:

\begin{lstlisting}
S [val=7, ops=2]
  x [lexval='x']
  = [lexval='=']
  E [val=7, ops=2]
    T [val=6, ops=1]
      F [val=2, ops=0]
        2 [lexval='2']
      M [inh=2, val=6, ops=1]
        * [lexval='*']
        F [val=3, ops=0]
          3 [lexval='3']
        M [inh=6, val=6, ops=0]
          e
    A [inh=6, val=7, ops=1]
      + [lexval='+']
      T [val=1, ops=0]
        F [val=1, ops=0]
          1 [lexval='1']
        M [inh=1, val=1, ops=0]
          e
      A [inh=7, val=7, ops=0]
        e
\end{lstlisting}

\subsubsection*{Grafo de dependencias}

El grafo de dependencias tiene una arista desde cada atributo usado hacia el atributo que define. Para \texttt{x=2*3+1}, las dependencias clave son:

\begin{itemize}[nosep]
  \item $2.\texttt{lexval} \to F.\texttt{val} \to M.\texttt{inh}$
  \item $3.\texttt{lexval} \to F'.\texttt{val}$; luego $M.\texttt{inh} \times F'.\texttt{val} \to M'.\texttt{inh} = 6$
  \item $M'.\texttt{val} = M'.\texttt{inh} = 6 \to M.\texttt{val} \to T.\texttt{val}$
  \item $T.\texttt{val} \to A.\texttt{inh} = 6$
  \item $1.\texttt{lexval} \to F''.\texttt{val} \to M''.\texttt{inh} \to M''.\texttt{val} \to T'.\texttt{val}$
  \item $A.\texttt{inh} + T'.\texttt{val} \to A'.\texttt{inh} = 7 \to A'.\texttt{val} \to A.\texttt{val} \to E.\texttt{val} \to S.\texttt{val} = 7$
\end{itemize}

El grafo es ac\'iclico, por lo que el orden topol\'ogico funciona y la evaluaci\'on procede de izquierda a derecha, de abajo hacia arriba para atributos sintetizados y de arriba hacia abajo para los heredados.

\subsubsection*{Orden de evaluaci\'on}

El orden topol\'ogico producido por nuestro programa (algoritmo de Kahn, empates rotos por id de nodo):

\begin{enumerate}[nosep]
  \item $F.\texttt{val} = 2$ \quad (del lexval de $2$)
  \item $M.\texttt{inh} = 2$
  \item $F'.\texttt{val} = 3$ \quad (del lexval de $3$)
  \item $M'.\texttt{inh} = 2 \times 3 = 6$
  \item $M'.\texttt{val} = 6$; $M.\texttt{val} = 6$; $T.\texttt{val} = 6$
  \item $A.\texttt{inh} = 6$
  \item $F''.\texttt{val} = 1$; $M''.\texttt{inh} = 1$; $M''.\texttt{val} = 1$; $T'.\texttt{val} = 1$
  \item $A'.\texttt{inh} = 6 + 1 = 7$; $A'.\texttt{val} = 7$; $A.\texttt{val} = 7$
  \item $E.\texttt{val} = 7$; $S.\texttt{val} = 7$
\end{enumerate}

Resultado: $x = 7$, operaciones: 2 (una multiplicaci\'on, una suma).

\subsection{Pasos 5--6: Implementaci\'on y pruebas}

\subsubsection*{Implementaci\'on}

La SDD se almacena como datos en \texttt{sdd\_definition.py}: para cada producci\'on (indexada 0--10), una lista de reglas, donde cada regla declara qu\'e atributo define, qu\'e atributos usa (por posici\'on), y un lambda que calcula el valor. El evaluador en \texttt{sdd\_evaluator.py} es \emph{gen\'erico}: recibe la gram\'atica, el diccionario de reglas y las declaraciones de atributos, y no sabe nada sobre la SDD espec\'ifica.

El evaluador:
\begin{enumerate}[nosep]
  \item Parsea la entrada usando el parser LL(1) (Algoritmo~4.34) y construye un \'arbol de parseo (objetos \texttt{TreeNode}).
  \item Asigna ids \'unicos a los nodos (recorrido pre-order).
  \item Construye el grafo de dependencias instanciando las reglas de cada producci\'on en cada nodo interno.
  \item Calcula un orden topol\'ogico usando el algoritmo de Kahn. Si el orden es incompleto (no todos los v\'ertices son alcanzados), la SDD es circular.
  \item Eval\'ua los atributos en orden topol\'ogico.
\end{enumerate}

\subsubsection*{Resultados de las pruebas}

\begin{center}
\begin{tabular}{@{}llccl@{}}
\toprule
\textbf{Cadena} & \textbf{Estado} & \textbf{val} & \textbf{ops} & \textbf{Prop\'osito} \\
\midrule
\texttt{x=1}         & aceptada & 1 & 0 & cadena m\'as corta \\
\texttt{x=2+3}       & aceptada & 5 & 1 & suma simple \\
\texttt{x=2*3+1}     & aceptada & 7 & 2 & prueba de precedencia \\
\texttt{x=(1+2)*3}   & aceptada & 9 & 2 & anidamiento profundo \\
\texttt{x=1+2+3}     & aceptada & 6 & 2 & suma recursiva a derecha \\
\texttt{1+2}          & rechazada & --- & --- & no est\'a en el lenguaje (falta \texttt{x=}) \\
\texttt{+}            & rechazada & --- & --- & no est\'a en el lenguaje \\
$\varepsilon$         & rechazada & --- & --- & cadena vac\'ia \\
\bottomrule
\end{tabular}
\end{center}

Todos los resultados son matem\'aticamente correctos y coinciden con el c\'alculo a mano.

\subsubsection*{Detecci\'on de circularidad}

Para probar la detecci\'on de circularidad, modificamos temporalmente la regla de $A \to \varepsilon$: en vez de $A.\texttt{val} = A.\texttt{inh}$, pusimos $A.\texttt{val} = A.\texttt{val}$ (una auto-dependencia). Ejecutando con la bandera \texttt{--circular}:

\begin{lstlisting}
$ python sdd_main.py --circular
(Using circular SDD for testing)

x=1
The SDD is circular for this input.
\end{lstlisting}

El orden topol\'ogico detecta el ciclo y lo reporta correctamente.


%% ============================================================
\section{Problemas encontrados y soluciones}
\label{sec:problems}

\begin{enumerate}
  \item \textbf{Epsilon en conjuntos FIRST.}
    La primera versi\'on del ciclo de punto fijo de FIRST no propagaba $\varepsilon$ correctamente cuando un no-terminal derivaba $\varepsilon$ a trav\'es de una cadena (ej.\ $A \to B$, $B \to \varepsilon$).
    La soluci\'on fue verificar si todos los s\'imbolos en el cuerpo de una producci\'on pueden derivar $\varepsilon$ antes de agregar $\varepsilon$ al conjunto de la cabeza.

  \item \textbf{Conflictos SLR(1) con cuerpos epsilon.}
    Cuando una producci\'on como $B \to \varepsilon$ se representa con un cuerpo vac\'io, el \'item de reducci\'on tiene el punto en la posici\'on~0 y \texttt{is\_reduce()} retornaba verdadero inmediatamente.
    Tuvimos que asegurarnos de que \texttt{\_slr\_body} mapee \texttt{[e]} a una tupla vac\'ia para que la maquinaria de \'items LR lo maneje correctamente, y que FOLLOW se use para las acciones de reducci\'on.
\end{enumerate}


%% ============================================================
\section*{Referencias}

Aho, A.~V., Lam, M.~S., Sethi, R., \& Ullman, J.~D. (2007).
\emph{Compilers: Principles, Techniques, and Tools} (2da ed.).
Pearson / Addison-Wesley.

\end{document}
